% ================================
% Modelo de Relatório Técnico – Fundep / Fiotec
% Compilar com: pdflatex (duas vezes)
% ================================
\documentclass[12pt,a4paper]{article}

% --- Pacotes básicos
\usepackage[utf8]{inputenc}
\usepackage[T1]{fontenc}
\usepackage[brazil]{babel}
\usepackage{geometry}
\usepackage{setspace}
\usepackage{graphicx}
\usepackage{array}
\usepackage{enumitem}
\usepackage{fancyhdr}
\usepackage{lastpage}
\usepackage{hyperref}
\usepackage{titlesec}
\usepackage{tabularx}
\usepackage{float}
\usepackage{array}

\graphicspath{{img/}{./img/}}

\newcolumntype{Y}{>{\raggedright\arraybackslash}X}
\newcolumntype{C}[1]{>{\centering\arraybackslash}p{#1}}
\newcolumntype{Z}{>{\centering\arraybackslash}X}

% --- Layout
\geometry{margin=2.5cm}
\onehalfspacing
\setlist[itemize]{left=0pt,itemsep=4pt,topsep=4pt}
\hypersetup{
  colorlinks=true,
  linkcolor=black,
  urlcolor=blue
}

% --- Cabeçalho/Rodapé
\pagestyle{fancy}
\fancyhf{}
\lhead{\textit{\ReportProject}}
\rhead{\textit{\ReportTitle}}
\cfoot{\thepage/\pageref*{LastPage}}
\setlength{\headheight}{15pt}

% --- Títulos mais compactos
\titlespacing*{\section}{0pt}{10pt}{6pt}
\titlespacing*{\subsection}{0pt}{8pt}{4pt}

% --- Campos parametrizados
\newcommand{\ReportTitle}{Planejamento Estratégico 2026-2030}
\newcommand{\ReportProject}{Tecnologia da Informação}
\newcommand{\ReportOrg}{Fundação de Desenvolvimento da Pesquisa}
\newcommand{\ReportAuthor}{André Vinícius de Oliveira}
\newcommand{\ReportRole}{Gerente de TI}
\newcommand{\ReportPeriod}{\textit{2026}}
\newcommand{\ReportVersion}{v0.1}
\newcommand{\ReportDate}{\today}

\begin{document}

% --- Capa
\begin{titlepage}
  \centering
  % (Opcional) insira logos quando quiser:
  % \includegraphics[height=1.4cm]{logo-fundep.png}\hfill
  % \includegraphics[height=1.4cm]{logo-fiotec.png}

  \vspace*{3cm}
  {\Large \textbf{\ReportOrg}\par}
  \vspace{0.8cm}
  {\huge \textbf{\ReportTitle}\par}
  \vspace{0.4cm}
  {\Large \textbf{\ReportProject}\par}
  \vspace{0.8cm}
  {\large Período de referência: \ReportPeriod \par}
  \vspace{0.8cm}
  \begin{tabular}{>{\bfseries}p{3.2cm}p{10.5cm}}
    Elaborado por: & \ReportAuthor \\
    Função:        & \ReportRole \\
    Versão:        & \ReportVersion \\
    Data:          & \ReportDate \\
  \end{tabular}

  \vfill
  \vspace*{1cm}
\end{titlepage}

% --- Sumário
\tableofcontents
\newpage

% --- Introdução
\section{Introdução}

O presente documento apresenta o Planejamento Estratégico de Tecnologia da Informação para o período 2026-2030, elaborado em resposta à solicitação da Diretoria da Fundação de Desenvolvimento da Pesquisa (Fundep). Este planejamento busca alinhar as iniciativas tecnológicas com os objetivos organizacionais, garantindo a modernização e a eficiência operacional da instituição.

\subsection{Contexto e Justificativa}

A transformação digital é essencial para o fortalecimento da capacidade institucional da Fundep. O setor de Tecnologia da Informação assume papel central nessa transformação, promovendo modernização de infraestrutura, otimização de processos e geração de valor aos stakeholders. Este planejamento estabelece as diretrizes e ações estratégicas que nortearão os investimentos e iniciativas de TI nos próximos cinco anos.

\subsection{Objetivos do Planejamento}

Este documento tem como objetivos:

\begin{itemize}
  \item Definir a visão e estratégia de tecnologia para a Fundep no período 2026-2030;
  \item Identificar projetos prioritários de modernização e transformação digital;
  \item Estabelecer indicadores de desempenho (KPIs) para acompanhamento de resultados;
  \item Alinhar os investimentos tecnológicos com as prioridades institutionais;
  \item Garantir a sustentabilidade e a evolução contínua da infraestrutura de TI.
\end{itemize}

\subsection{Escopo e Abrangência}

O escopo deste planejamento abrange toda a área de Tecnologia da Informação da Fundep, incluindo infraestrutura de sistemas, operações, segurança da informação, modernização de aplicações e suporte ao usuário. O horizonte de planejamento estende-se de 2026 a 2030, com revisões anuais para ajuste de prioridades conforme a evolução organizacional.

\end{document}